%% evaluation.tex
%%

%% ==================
\section{Evaluation}
\label{sec:Evaluation}
%% ==================

This section presents the experimental results from evaluating all three scaling policies (Reactive HPA, Aggressive PAKS, Stability-Aware PAKS) across 5 independent random seeds (42--46). Each seed generates a different synthetic workload realization and trains a fresh predictor, so results represent performance averaged over different demand patterns.

\subsection{Overall Results}

Table~\ref{tab:results} reports mean $\pm$ standard deviation across all 5 seeds for the four evaluation metrics defined in Section~\ref{sec:Methodology}.

\begin{table}[htp]
\centering
\caption{Autoscaling performance across 5 seeds (mean $\pm$ std)}
\label{tab:results}
\begin{tabular}{lcccc}
\hline
\textbf{Model} & \textbf{SLA Violations} & \textbf{Over-prov. \%} & \textbf{Scaling Events} & \textbf{Pod Volatility} \\
\hline
Reactive HPA & 20.0 $\pm$ 2.7 & 40.92 $\pm$ 0.19 & 407.4 $\pm$ 4.9 & 3.32 $\pm$ 0.25 \\
Aggressive PAKS & 11.0 $\pm$ 2.2 & 45.80 $\pm$ 2.23 & 384.4 $\pm$ 9.1 & 2.86 $\pm$ 0.36 \\
Stability-Aware PAKS & 18.0 $\pm$ 2.7 & 46.52 $\pm$ 2.16 & 165.8 $\pm$ 5.3 & 1.95 $\pm$ 0.17 \\
\hline
\end{tabular}
\end{table}

\paragraph{Interpretation Guide}
\begin{itemize}
  \item \textbf{SLA Violations}: Lower is better. Measures how often the system under-provisions, causing service degradation.
  \item \textbf{Over-provisioning \%}: Lower is better. Quantifies resource waste and cost.
  \item \textbf{Scaling Events}: Lower is better. High event counts indicate instability and operational overhead (each scaling action incurs startup latency, configuration overhead, and potential service disruption).
  \item \textbf{Pod Volatility}: Lower is better. Captures both frequency and magnitude of scaling actions; penalizes large swings more than small adjustments.
\end{itemize}

The following subsections analyze each comparison in detail.

%% ===============================
\subsection{RQ1: Reproducing NimbusGuard's Agility-Instability Trade-off}
\label{sec:evalOne}
%% ===============================

\textbf{Research Question}: Does a simple proactive scaler (Aggressive PAKS) reproduce the trade-off NimbusGuard reports: better SLA compliance than reactive HPA, but at a cost in resource efficiency and/or stability?

\textbf{Comparison}: Aggressive PAKS vs. Reactive HPA (both baselines).

\paragraph{SLA Violations} Aggressive PAKS achieves 11.0 $\pm$ 2.2 violations compared to HPA's 20.0 $\pm$ 2.7, a 45\% reduction. The difference (9.0 violations) is more than 3× the larger standard deviation, indicating a statistically meaningful improvement. This confirms that proactive forecasting enables the system to scale up \emph{before} spikes arrive, avoiding the under-provisioning that HPA suffers due to its one-step reaction delay.

\paragraph{Over-provisioning} Aggressive PAKS incurs 45.80\% $\pm$ 2.23\% over-provisioning compared to HPA's 40.92\% $\pm$ 0.19\%, a 4.88 percentage-point increase. This represents approximately 12\% higher resource waste. The cost is attributable to two factors:
\begin{enumerate}
  \item The 5\% safety buffer applied to predictions (by design, to guard against under-prediction).
  \item Prediction errors: when the predictor overestimates future workload, Aggressive PAKS scales up more than necessary.
\end{enumerate}

\paragraph{Scaling Events} Interestingly, Aggressive PAKS (384.4 $\pm$ 9.1 events) and Reactive HPA (407.4 $\pm$ 4.9 events) have \emph{comparable} scaling-event counts, differing by only 23 events (5.6\%). This contrasts with NimbusGuard's reported 2× gap (8 scaling events for NimbusGuard vs. 4 for HPA/KEDA). The difference in specific numbers is expected -- this is a much simpler feed-forward predictor on a synthetic workload rather than a DQN+LSTM agent on a real Kubernetes testbed -- but the \emph{qualitative} finding holds: Aggressive PAKS is not dramatically more unstable than HPA in event count, despite acting on predictions.

\paragraph{Pod Volatility} Aggressive PAKS achieves lower volatility (2.86 $\pm$ 0.36) than Reactive HPA (3.32 $\pm$ 0.25), a 13.9\% reduction. This suggests that Aggressive PAKS's scaling actions, while frequent, are \emph{smoother} than HPA's reactive adjustments. HPA's one-step delay causes it to overshoot (scaling up too much after a spike ends, then scaling back down), leading to larger step-to-step changes.

\paragraph{Answer to RQ1} \textbf{Yes.} Aggressive PAKS reproduces the core trade-off: significantly better SLA compliance than reactive HPA (45\% fewer violations), at the cost of moderately higher over-provisioning (12\% increase). The scaling-event count does not exhibit the 2× instability gap NimbusGuard reports, likely because this simpler predictor on a synthetic workload produces smoother predictions than a DQN agent would. Nonetheless, the motivation for stability improvement remains: over-provisioning is higher, and a more aggressive predictor (or noisier workload) could easily push scaling-event counts higher.

%% ===============================
\subsection{RQ2: Can Stability Be Recovered Without Sacrificing SLA Performance?}
\label{sec:evalTwo}
%% ===============================

\textbf{Research Question}: Does the stability-aware controller (exponential smoothing + hysteresis/cooldown) substantially reduce instability metrics (scaling events, volatility) while retaining most of the SLA benefit over reactive HPA?

\textbf{Comparison}: Stability-Aware PAKS vs. Aggressive PAKS (both proactive).

\paragraph{Scaling Events} Stability-Aware PAKS achieves 165.8 $\pm$ 5.3 events compared to Aggressive PAKS's 384.4 $\pm$ 9.1 events, a \textbf{57\% reduction} (218.6 fewer events). The difference is 24× the larger standard deviation, indicating an extremely robust improvement. This directly validates the core hypothesis: smoothing and hysteresis/cooldown dramatically reduce the frequency of scaling actions.

Figure~\ref{fig:scaling_events_timeseries} visualizes this effect over a representative 100-timestep window, showing that Aggressive PAKS scales almost every few steps in response to prediction fluctuations, while Stability-Aware PAKS scales only when sustained trends justify action.

\paragraph{Pod Volatility} Stability-Aware PAKS achieves 1.95 $\pm$ 0.17 volatility compared to Aggressive PAKS's 2.86 $\pm$ 0.36, a \textbf{32\% reduction}. The difference (0.91) is 2.5× the larger standard deviation. This confirms that not only does Stability-Aware PAKS scale \emph{less often}, but when it does scale, the changes are \emph{smaller on average} (because hysteresis prevents reaction to small differences, so only larger sustained shifts trigger action).

\paragraph{SLA Violations} Stability-Aware PAKS incurs 18.0 $\pm$ 2.7 violations compared to Aggressive PAKS's 11.0 $\pm$ 2.2 violations, a 63.6\% \emph{increase} (7 additional violations). The difference (7.0) is 2.6× the larger standard deviation, so this is a real cost, not noise. However, Stability-Aware PAKS still beats Reactive HPA (18.0 vs. 20.0), a 10\% improvement.

The SLA cost arises because smoothing and cooldown necessarily delay reaction to genuine demand shifts, not just noise. When a real spike arrives, Aggressive PAKS can respond immediately (potentially on the very next timestep), while Stability-Aware PAKS must wait for:
\begin{enumerate}
  \item The EMA to accumulate the signal (multiple predictions pointing in the same direction).
  \item The hysteresis threshold to be exceeded (change must be $>1$ pod).
  \item The cooldown timer to expire (if a recent scaling action occurred).
\end{enumerate}
During this delay, the system may under-provision for a few timesteps.

\paragraph{Over-provisioning} Stability-Aware PAKS (46.52\% $\pm$ 2.16\%) and Aggressive PAKS (45.80\% $\pm$ 2.23\%) have essentially \emph{identical} over-provisioning, differing by only 0.72 percentage points (1.6\% relative difference, well within one standard deviation). This indicates that the stability fix targets \emph{when} scaling actions occur, not the long-run resource allocation level. Both policies use the same predictor and the same 5\% safety buffer, so they converge to similar average allocations over time; the difference is in the path taken to get there.

\paragraph{Answer to RQ2} \textbf{Partially.} Stability is recovered \emph{substantially}: 57\% fewer scaling events and 32\% lower volatility, addressing the specific instability weakness NimbusGuard's own results document. However, this comes at a cost: 7 additional SLA violations (a 63.6\% increase relative to Aggressive PAKS, though still better than Reactive HPA). The result is not a ``free lunch'' but a \emph{different point on the trade-off curve}: Stability-Aware PAKS sacrifices some of Aggressive PAKS's agility to gain stability, while still outperforming the reactive baseline.

%% ===============================
\subsection{RQ3: Which Policy Should Be Deployed?}
\label{sec:evalThree}
%% ===============================

\textbf{Research Question}: Given the measured trade-offs, which policy is ``best''?

\textbf{Answer}: \textbf{It depends on the operator's priorities.} This is not a single-winner scenario but a Pareto frontier with three distinct points:

\paragraph{Reactive HPA (Baseline 1)}
\begin{itemize}
  \item \textbf{Strengths}: Lowest over-provisioning (40.92\%), no forecasting complexity, well-understood behavior.
  \item \textbf{Weaknesses}: Highest SLA violations (20.0), highest volatility (3.32), high scaling-event count (407.4).
  \item \textbf{Use case}: Cost-sensitive deployments where occasional under-provisioning is tolerable and operational simplicity is valued.
\end{itemize}

\paragraph{Aggressive PAKS (Baseline 2)}
\begin{itemize}
  \item \textbf{Strengths}: Lowest SLA violations (11.0), lowest volatility (2.86), demonstrates proactive forecasting capability.
  \item \textbf{Weaknesses}: Highest over-provisioning (45.80\%), moderately high scaling-event count (384.4).
  \item \textbf{Use case}: Performance-critical deployments where SLA violations are expensive (e.g., revenue-generating APIs, real-time services) and the 12\% resource-cost increase is acceptable.
\end{itemize}

\paragraph{Stability-Aware PAKS (Proposed)}
\begin{itemize}
  \item \textbf{Strengths}: Dramatically lowest scaling-event count (165.8, 57\% reduction vs. Aggressive PAKS), lowest volatility (1.95), still beats HPA on SLA (18.0 vs. 20.0).
  \item \textbf{Weaknesses}: More SLA violations than Aggressive PAKS (18.0 vs. 11.0), over-provisioning comparable to Aggressive PAKS (46.52\%).
  \item \textbf{Use case}: Deployments where scaling actions themselves are costly or disruptive (e.g., stateful services with long startup times, systems with tight SLO budgets on reconfiguration overhead), and the operator can tolerate a few additional SLA violations in exchange for operational stability.
\end{itemize}

\paragraph{Decision Framework}
If the operator assigns costs to each metric (e.g., \$X per SLA violation, \$Y per percentage-point of over-provisioning, \$Z per scaling event), the ``best'' policy is the one minimizing total cost:
\[
\text{Cost} = C_{\text{SLA}} \cdot \text{violations} + C_{\text{over-prov}} \cdot \text{over-prov\%} + C_{\text{event}} \cdot \text{events} + C_{\text{volatility}} \cdot \text{volatility}
\]

In the absence of such costs, the honest answer is that this project documents the trade-off surface, not a single winner. Section~\ref{sec:Conclusion} discusses future work on tuning the smoothing/hysteresis/cooldown parameters to explore the full range of this trade-off.

%% ===============================
\subsection{Variability Across Seeds}
\label{sec:variability}
%% ===============================

Standard deviations are relatively small for most metrics (e.g., over-provisioning varies by $<3$ percentage points, scaling events by $<10$ events), indicating that findings are robust across different workload realizations. The largest relative variability is in SLA violations (std $\approx$ 13--25\% of mean), which is expected: under some workload realizations, the specific timing of spikes relative to the predictor's lookback window may make violations more or less likely. Nonetheless, the rank-ordering of policies (Aggressive PAKS best on SLA, Stability-Aware PAKS best on events/volatility) holds consistently across all seeds.

%% ===============================
\subsection{Comparison to NimbusGuard (Baseline Paper)}
\label{sec:comparison_nimbusguard}
%% ===============================

Direct numerical comparison to NimbusGuard is inappropriate because:
\begin{enumerate}
  \item NimbusGuard evaluates on a real Kubernetes testbed with a specific phased load pattern; this project uses a synthetic workload.
  \item NimbusGuard uses a DQN+LSTM+LLM pipeline; this project uses a feed-forward MLPRegressor.
  \item NimbusGuard reports ``Total Scaling Events'' as an integer count over one experiment run; this project averages over 5 seeds with standard deviations.
\end{enumerate}

However, the \emph{qualitative} findings align:
\begin{itemize}
  \item Proactive scaling beats reactive baselines on SLA/responsiveness metrics (\checkmark reproduced).
  \item Proactive scaling incurs higher resource costs (\checkmark reproduced: over-provisioning increases by 12\%).
  \item Unfiltered proactive scaling can exhibit instability (\checkmark reproduced: Aggressive PAKS has high scaling-event counts and volatility).
  \item Stability mechanisms can mitigate this instability at a cost (\checkmark demonstrated: Stability-Aware PAKS cuts events by 57\% but adds 7 SLA violations).
\end{itemize}

The contribution is not in matching NimbusGuard's exact numbers, but in \emph{reproducing the trade-off structure} with a simpler, fully-specified approach and \emph{building the stability fix NimbusGuard names as future work}.

%% ===============================
\subsection{Threats to Validity}
\label{sec:threats}
%% ===============================

\paragraph{Internal Validity} All comparisons are controlled: the same workload per seed is used for all three policies, and the same predictor is shared between Aggressive and Stability-Aware PAKS. Results are deterministic given the random seed. The primary threat is parameter choice -- the specific values for $\alpha$, hysteresis threshold, and cooldown were chosen based on literature precedent and limited empirical observation on a single seed, not exhaustive tuning. Different parameter configurations could yield different points on the trade-off curve.

\paragraph{External Validity} The synthetic workload may not capture all characteristics of real cloud workloads (e.g., long-tail distributions, correlated arrivals across services, non-stationary patterns). Evaluation on real traces (Google Cluster, Alibaba, Azure Functions) and on a real Kubernetes testbed (as NimbusGuard does) would strengthen generalizability.

\paragraph{Construct Validity} The four metrics (SLA violations, over-provisioning, scaling events, volatility) are standard in the autoscaling literature and directly comparable to those NimbusGuard reports. However, they do not capture all costs -- for example, the latency overhead of starting/stopping pods, or the impact of pod churn on service quality (e.g., connection draining). A real deployment would need to measure end-to-end latency and availability.

\paragraph{Conclusion Validity} Results are averaged over 5 seeds and standard deviations are reported. Differences deemed ``meaningful'' exceed one standard deviation. For stronger statistical confidence, more seeds (e.g., 10--20) would be ideal, though the current sample size is sufficient to demonstrate the core trade-offs.
